\title{Parameter-Efficient Orthogonal Finetuning via Factorization}

\begin{abstract}
The widespread adoption of large foundation models is limited by the immense computational demands associated with training such models from the ground up. Consequently, methods that facilitate the efficient adaptation of these models for specific downstream tasks are becoming ever more critical. This study investigates Orthogonal Finetuning (OFT), a theoretically grounded approach for tailoring foundation models to downstream applications. While OFT offers robust generalization, it requires a considerable number of trainable parameters due to the intrinsic high dimensionality of orthogonal matrices. To overcome this challenge, we analyze OFT through the lens of information transfer and distill several key principles that promote greater parameter efficiency. Drawing inspiration from the Cooley-Tukey fast Fourier transform, renowned for its efficiency in transmitting information, we design a parameterization based on butterfly architectures. Integrating this structure with OFT yields a new method—Orthogonal Butterfly (BOFT)—that significantly improves parameter efficiency. BOFT generalizes OFT, treating it as a particular instance within a broader orthogonal finetuning framework. We empirically validate BOFT’s effectiveness by extensively adapting large-scale vision transformers, language models, and text-to-image diffusion models to a variety of vision and language tasks.
\end{abstract}

\section{Introduction}

State-of-the-art systems such as ChatGPT~\cite{brown2020language,chen2021evaluating} and Stable Diffusion~\cite{rombach2022high} showcase the extraordinary generalization capacity of modern large foundation models. The remarkable improvement in their capabilities, however, has been coupled with a surge in parameter count—

(\emph{e.g.}, the GPT-3 architecture comprises nearly 175B parameters). This substantial growth makes it infeasible for most researchers to train such models anew. As a result, advancing the field now hinges on the ability to efficiently repurpose pre-trained foundation models for specialized tasks without full retraining. To address this, three principal adaptation strategies have emerged: \emph{model finetuning}~\cite{hulora2022,zaken2022bitfit,guo2020parameter,raffel2020exploring,qiu2023controlling,zhang2022adaptive,chavan2023one}, which leaves the model’s structure intact but updates only select parameters; \emph{adapter tuning}~\cite{rebuffi2017learning,houlsby2019parameter,pfeiffer2020adapterfusion,lin2020exploring,he2021towards}, which incorporates trainable adapter modules into the existing architecture; and \emph{prompt tuning}~\cite{li2021prefix,lester2021power}, which introduces tunable prefix tokens as part of the model’s input. Among these, model finetuning stands out due to its straightforward implementation, powerful adaptation capability, and the critical advantage of not introducing any extra inference latency.

The core idea of model finetuning is to maintain similarity between the original pretrained model and its finetuned version according to certain metrics, thereby safeguarding the knowledge acquired during pretraining. A common tactic in existing finetuning strategies is the use of a reduced learning rate, an informal practice that naturally restricts changes and hence the divergence between the initial and the adapted models. Given that weight matrices have a structured composition, a more systematic strategy involves retaining the relationships embedded within these matrices, specifically, the pairwise angular relations between neurons. This realization has led to the development of the Orthogonal Finetuning (OFT) framework~\cite{qiu2023controlling}, where an orthogonal matrix is uniformly applied to all neurons within the same layer, ensuring that the inter-neuronal angles remain invariant during finetuning. OFT has been shown to provide both robust generalization and strong convergence, particularly when finetuning text-to-image diffusion models~\cite{qiu2023controlling}. Nevertheless, due to the large size of orthogonal matrices in high-dimensional settings, OFT often involves a significant number of parameters. To mitigate this issue, the framework employs a block-diagonal configuration, effectively reducing the parameter count. This enhanced efficiency, however, imposes certain limitations: the imposed block-diagonal pattern fixes the sparsity structure of the orthogonal matrix, and each block operates its transformation independently. While this parameter-saving approach does alleviate computational costs, it may introduce unfavorable inductive biases (such as reduced model expressiveness and an inability to represent classical linear mappings), and crucially, the optimal selection of the sparsity pattern remains unresolved.

The central challenge lies in constructing a dense orthogonal matrix in a way that conserves parameters. Although creating a $d$-dimensional dense orthogonal matrix conventionally requires $\mathcal{O}(d^2)$ parameters—making it seemingly impractical for large $d$—we introduce an alternative approach: assembling a dense orthogonal matrix from a sequence of factorized sparse matrices. This strategy leverages the principle that computational workload can be exchanged for reduced parameter usage. Because our representation composes the orthogonal matrix as a product of sparse matrices, each training iteration necessitates repeatedly multiplying these sparse factors. To contextualize this matrix factorization, we recast the task of generating a dense orthogonal matrix as an information transmission challenge. Under this interpretation, the multiplication of sparse matrices is analogous to transmitting data across a grid-shaped network. This perspective opens up multiple potential sparse factorization schemes that constrain the number of trainable parameters without sacrificing the expressiveness required to produce dense matrices.

For parameter efficiency within this information transmission paradigm, we draw upon the butterfly structure from the Cooley-Tukey fast Fourier transform algorithm~\cite{cooley1965algorithm}, which enables highly efficient data exchange among nodes~\cite{klasing1994broadcasting}. Notably, the butterfly graph structure within the Cooley-Tukey framework naturally gives rise to a sparse factorization method termed butterfly factorization. Utilizing butterfly factorization, one can represent a $d$-dimensional dense matrix as a product of $\mathcal{O}(\log d)$ sparse matrices, each containing only $\mathcal{O}(d)$ nonzero entries. By further ensuring the orthogonality of these sparse matrices, we obtain an efficient orthogonal parameterization that requires just $\mathcal{O}(d\log d)$ parameters—representing a dramatic reduction compared to the full parameter count. Building upon this efficient orthogonal structure, we introduce a novel, parameter-saving finetuning method: Orthogonal Butterfly~(BOFT). BOFT generalizes and extends OFT as a special case, establishing a comprehensive framework for orthogonal finetuning. Both block-diagonal and butterfly constructions share a core property: they impose sparsity on the orthogonal matrices to diminish the number of trainable parameters. It is instructive to compare our technique with the low-rank methods exemplified by LoRA~\cite{hulora2022}; both low-rank and sparse matrices constitute key categories of structured matrices~\cite{candes2011robust} that facilitate efficient parameter usage.

In contrast to OFT, which leverages a block-diagonal architecture to balance model expressivity with structural regularity, BOFT adopts the butterfly architecture to facilitate a more seamless transition between elements of the full orthogonal group (providing expressiveness) and the identity (ensuring regularity). This approach allows for the identification of a more compact hypothesis space within the orthogonal group that better serves subsequent learning objectives. The butterfly structure, being foundational in numerous efficient linear operators such as the discrete Fourier and discrete cosine transforms, motivates our argument that embedding such a structure in parameter-efficient modeling brings a useful inductive bias. This bias, in turn, promotes better generalization and offers resistance to overfitting. Our reasoning is based on the observation that the butterfly structure can flexibly approximate a range of classic linear operations, a versatility lacking in the block-diagonal configurations utilized in OFT. The main contributions of our work are summarized as follows:

\begin{itemize}[leftmargin=*,nosep]
    \item We address parameter-efficient orthogonal finetuning by introducing a novel framework rooted in information transmission, reformulating the challenge of constructing a dense, parameter-efficient orthogonal matrix as a problem of information flow across a graph with grid topology.
    \item Building on the butterfly architecture employed in the Cooley-Tukey algorithm, we introduce Orthogonal Butterfly, a new method for orthogonal finetuning, situated within our information transmission framework.
    \item We offer theoretical perspectives that help elucidate why BOFT is able to drastically reduce the count of learnable parameters without compromising expressivity or stability during optimization. Owing to the matrix factorization properties, BOFT also enables an intuitive interpolation of model weights (refer to Figure~\ref{fig:boft_interpolation}).
    \item For the first time, we extend orthogonal finetuning methodologies beyond their prior applications in controllable text-to-image generation~\cite{qiu2023controlling}, demonstrating the robustness and versatility of our method as a general finetuning strategy. We test BOFT on a variety of adaptation tasks covering domains from computer vision to natural language processing, where it not only exceeds contemporary methods in performance but also exemplifies remarkable parameter efficiency and generalization prowess.
\end{itemize}

\section{Related Work}

\textbf{Parameter-efficient finetuning (PEFT)}. As foundation models (\emph{e.g.}, \cite{brown2020language,radford2021learning,rombach2022high,kirillov2023segment}) continue to increase in size and capacity, the development of parameter-efficient approaches for adapting these models to downstream applications has garnered significant attention~\cite{treviso2023efficient,ding2023parameter,lialin2023scaling}. Among the array of PEFT strategies~\cite{houlsby2019parameter,aghajanyan2020intrinsic,hulora2022,edalati2022krona,wang2022adamix,gheini2021cross,zaken2022bitfit,guo2020parameter,sung2021training,ansell2022composable,lester2021power,li2021prefix,vu2022spot,he2021towards,mao2021unipelt,karimi2021compacter,liu2022few,sung2022lst,chen2022parameter,jia2022visual,chen2022adaptformer,zhang2022neural,jie2023fact,lian2022scaling,luo2023towards}, those based on reparameterization~\cite{aghajanyan2020intrinsic,hulora2022,edalati2022krona} are particularly pertinent for our purposes. The LoRA technique~\cite{hulora2022} augments pretrained weight matrices through the addition of products of two low-rank matrices, achieving notable results on various natural language processing tasks. Whereas LoRA assigns a fixed rank to every added matrix, recent works \cite{zhang2022adaptive,valipour2022dylora,zhang2023increlora} explore dynamically varying the rank at each layer to better distribute the available parameter resources. Because of the straightforwardness of low-rank reparameterization, this approach has seen widespread adoption \cite{dettmers2023qlora,zi2023delta,chavan2023one}. Drawing on the insight that hyperspherical energy is associated with generalization~\cite{liu2018learning,liu2021orthogonal}, \cite{qiu2023controlling} introduce orthogonal finetuning (OFT), a compelling alternative for finetuning text-to-image diffusion models. In OFT, an orthogonal matrix is trained to perform intra-layer neuron transformations, resulting in improved generalization and more robust training as compared to LoRA. However, this gain comes at the cost of generally increased trainable parameter counts in OFT relative to LoRA. Thus, optimizing OFT for greater parameter efficiency remains a meaningful direction. Furthermore, it remains an open question whether the advantages of OFT extend beyond text-to-image diffusion model adaptation~\cite{qiu2023controlling}. BOFT addresses the parameter efficiency limitation of OFT through the use of butterfly factorization, enabling us to investigate the benefits of orthogonal finetuning on a broader range of adaptation problems.

\textbf{Butterfly structures}. The radix-2 Cooley-Tukey algorithm achieves an efficient computation of the $N$-point discrete Fourier transform by recursively decomposing it into two $\frac{N}{2}$-point DFTs, inherently generating a butterfly pattern representable as a product of sparse matrices (collectively known as a butterfly matrix). Such matrices have been utilized to parameterize orthogonal transformations—helping to eliminate pivoting during Gaussian elimination and boosting computational efficiency~\cite{parker1995random}—and to promote stable training in recurrent neural networks~\cite{jing2017tunable}, as well as in kernel approximation~\cite{munkhoeva2018quadrature}. Fast linear transforms have been learned via butterfly-based parameterizations in~\cite{dao2019learning,dao2019kaleidoscope}, while \cite{chen2022pixelated,dao2022monarch} have employed butterfly matrices to facilitate scalable neural network training. Further, butterfly constructions are effective for rapid matrix-vector products~\cite{michielssen1996multilevel,o2010algorithm}, compact matrix approximations~\cite{li2015butterfly}, and network transmission~\cite{klasing1994broadcasting,li2003linear,rajasingh2016transmission}. Unlike these prior usages, our research focuses on leveraging butterfly structures specifically to make OFT more parameter efficient.

\section{An Information Transmission View on Orthogonal Finetuning}

We begin by reviewing foundational concepts related to OFT. In OFT, the adaptation of a pretrained weight matrix involves reparameterizing the new matrix as the product of a trainable orthogonal matrix and the fixed pretrained weights. While LoRA modifies weights through an additive low-rank component, OFT instead updates the weights via a multiplicative orthogonal matrix. For parameter-efficiency, LoRA relies on a low-rank formulation; in contrast, the original OFT adopts a block-diagonal orthogonal approach~\cite{qiu2023controlling}, where reducing the diagonal block size increases efficiency. Figure~\ref{fig:comp_lora} visually contrasts these methods. The rationale for finetuning with an orthogonal transformation is to maintain neuron pairwise angles~\cite{liu2018learning,liu2021orthogonal,qiu2023controlling}, thereby largely retaining the semantic content acquired during pretraining. More specifically, OFT learns an orthogonal matrix $\bm{R}\in\mathbb{R}^{d\times d}$ in the context of a pretrained linear layer $\bm{W}^0\in\mathbb{R}^{d\times n}$, altering the forward computation from $\bm{z}=(\bm{W}^0)^\top\bm{x}$ to $\bm{z}=(\bm{R}\bm{W}^0)^\top\bm{x}$, where $\bm{x}\in\mathbb{R}^{d}$ is the input and $\bm{z}\in\mathbb{R}^{n}$ the output. For zero initialization, $\bm{R}$ is set to the identity matrix. To maintain orthogonality of $\bm{R}$ during finetuning, Cayley parameterization is used as in~\cite{qiu2023controlling,liu2021orthogonal}, i.e., $\bm{R}=(\bm{I}+\bm{Q})(\bm{I}-\bm{Q})^{-1}$, where $\bm{Q}$ is a skew-symmetric matrix, satisfying $\bm{Q}=-\bm{Q}^\top$. Parameter savings are achieved by structuring $\bm{R}$ as block-diagonal, i.e., $\bm{R} = \text{diag}(\bm{R}_1, \bm{R}_2, \dots, \bm{R}_r)$, with each block $\bm{R}_i\in\mathbb{R}^{b\times b}$ (and $br=d$) being a compact orthogonal matrix. However, this block-diagonal design enforces an implicit assumption: each neuron's vector (a column of $\bm{W}^0$) is partitioned into $r$ separate groups, with each group being transformed independently by different orthogonal blocks. Although this scheme is empirically successful, the assignment of neuron dimensions to groups according to their indices is conceptually arbitrary, highlighting the potential benefits of using dense orthogonal matrices instead. This consideration naturally raises the question: \emph{Is it possible to design a dense orthogonal matrix while preserving parameter-efficiency?}

To tackle this problem, we suggest constructing a dense orthogonal matrix by multiplying several sparse orthogonal matrices together. To develop an integrated yet accessible understanding of how sparsity emerges in orthogonal matrix factorization, we recast the challenge of producing a dense orthogonal matrix in OFT as analogous to an information transmission task. More concretely, forming a dense matrix $\bm{R}\in\mathbb{R}^{d\times d}$ from a sequence of $m$ square matrices as $\thickmuskip=2mu \medmuskip=2mu \bm{R}=\bm{B}_m\bm{B}_{m-1}\cdots\bm{B}_1$ can be interpreted as relaying information through a grid structure comprising $d\times (m+1)$ nodes, as depicted in Figure~\ref{fig:info_trans}. This perspective is motivated by the insight that a $d$-by-$d$ dense matrix naturally reflects dense connectivity linking two sets of $d$ nodes. For instance, in matrix $\bm{R}$, a zero at entry $R_{ij}$ signifies an absence of information flow from node $j$ to node $i$, while a nonzero value permits transmission. Consequently, expressing $\bm{R}$ as the product $\bm{B}_m\bm{B}_{m-1}\cdots\bm{B}_1$ can be viewed as orchestrating a sequence of information exchanges governed by the graphs induced by each $\bm{B}_i$. The flow of information starts following the structure of $\bm{B}_1$ and culminates in $\bm{B}_n$. As an illustrative case in Figure~\ref{fig:info_trans}, we examine the decomposition $\bm{R}=\bm{B}_5\bm{B}_4\bm{B}_3\bm{B}_2\bm{B}_1$, where both the sparsity structures and the resulting graph are displayed. The graph in the figure arises by unraveling the steps in the matrix product. Here, $\bm{B}_i$ defines the connections between the $i$-th and $(i+1)$-th node layers. To be more precise, the $(j_1,j_2)$ entry in $\bm{B}_i$ specifies whether a directed connection runs from node $j_2$ in layer $i$ to node $j_1$ in layer $i+1$ (a zero entry indicates absence). For the product $\bm{B}_5\bm{B}_4\bm{B}_3\bm{B}_2\bm{B}_1$ to yield a fully dense matrix, every node at the first level must be able to send information to every node at the sixth level. Restricting attention to $\bm{R}=\bm{B}_4\bm{B}_3\bm{B}_2\bm{B}_1$, which encompasses communication from initial nodes (level one) to final nodes (level five), we observe that node $1$ fails to reach node $3$; thus, the $(3,1)$ entry is zero in the combined sparsity pattern of $\bm{B}_4\bm{B}_3\bm{B}_2\bm{B}_1$. The same analogy holds when further limiting the product to $\bm{R}=\bm{B}_3\bm{B}_2\bm{B}_1$ or $\bm{R}=\bm{B}_2\bm{B}_1$. In general, to ensure that a matrix $\bm{R}\in\mathbb{R}^{d\times d}$ is dense, the collection of $m$ matrices $\bm{B}_m,\cdots,\bm{B}_1$ must collectively establish a set of directed edges over a $d\times (m+1)$ lattice, with each directed edge joining only pairs of adjacent layers (i.e., columns), such that information originating from any node at the first level can ultimately be delivered to every node at the $(m+1)$-th level.

Adopting the information transmission framework offers valuable insights into the distribution of non-zero entries in the factorization matrices used in OFT. Figure~\ref{fig:blk_diag} provides an illustration of the block-diagonal configuration present in $\bm{R}$ within the original OFT setup. While this configuration effectively lowers the parameter count to be optimized, it is inherently incapable of forming a fully dense $\bm{R}$. Our objective is to generate a dense orthogonal matrix through the composition of $m$ sparse orthogonal matrices, aiming to use as few trainable parameters as feasible. Within the information transmission paradigm, achieving this goal requires two fundamental properties: \emph{(i)} \textbf{dense connectivity}, wherein each input node in the initial layer maintains at least one route to every output node in the terminal layer, and \emph{(ii)} \textbf{minimum free edges}, which mandates that the aggregate number of edges be minimized under orthogonality constraints. Notably, orthogonality introduces intricate restrictions on the permissible edge configurations between adjacent layers. In particular, ensuring that each $\bm{B}_i$ is full-rank—a precondition for orthogonality—demands that $d$ edges connect all nodes between the $i$-th and $(i=1)$-th levels, thereby setting a lower bound of $d$ edges for such pairs (\emph{e.g.}, 4 in Figure~\ref{fig:info_trans}). These required $d$ edges are determined solely by orthogonality and do not contribute to the tally of trainable edges, as their corresponding elements are fixed and not learnable (\emph{e.g.}, in a $d\times d$ orthogonal matrix with $d$ non-zeros, each entry must be $\pm 1$). Due to the fewer parameters required by orthogonal matrices in comparison to general matrices, this orthogonality constraint introduces additional dependencies among edges. As an illustrative case, in a $2\times 2$ orthogonal matrix, imposing a zero at the $(1,1)$ position necessitates a zero at the $(2,2)$ entry as well (\emph{i.e.}, the absence of one edge can entail the absence of another). Consequently, for a given edge configuration, the orthogonality condition may cause some edges to be either mandated or excluded. By examining the support pattern (i.e., non-zero structure) of the constructed orthogonal matrix, the information transmission perspective proves especially beneficial for OFT, since our focus is solely on the arrangement of trainable non-zero elements in $\bm{R}$, irrespective of their actual numerical values.

A straightforward dense linkage between two hierarchical levels involves $\mathcal{O}(d^2)$ connections (specifically, a single dense orthogonal matrix), resulting in $d^2-d$ total edges (for example, when $d=4$, there are 12 edges). Figure~\ref{fig:info_trans} illustrates an alternative matrix decomposition that requires only $10$ edges in total—fewer than those needed by one dense orthogonal matrix. This approach offers a graphical viewpoint for analyzing the parameter-efficiency of OFT, allowing straightforward derivation of valid decompositions using this scheme. Our methodology draws on an efficient topology from the Cooley-Tukey algorithm, known as butterfly graphs~\cite{cooley1965algorithm}, which facilitate a dense and efficient linkage between $d$ source nodes and $d$ target nodes using just $\mathcal{O}(d\log d)$ edges. For instance, the network topology shown in Figure~\ref{fig:info_trans} achieves dense connectivity with $10$ edges, whereas the butterfly network can accomplish this with only $8$ edges. In the following, we describe how the butterfly structure enhances parameter efficiency.

\section{Orthogonal Parameterization by Butterfly Factorization}

The butterfly architecture was originally developed for use in the Cooley-Tukey algorithm to accelerate the fast Fourier transform. In the context of the Fourier transform, a minor modification in the frequency space produces a comprehensive effect in the spatial domain, paralleling our scenario for information propagation: each node at the initial level can transmit data to all nodes at the final level. Beyond signal processing, the butterfly structure is well-established as an effective topology in computer networks~\cite{solihin2015fundamentals,li2011linear} for facilitating efficient data transfer. Let $k \geq 2$ be a power of $2$. We begin by introducing the \emph{butterfly factor} $\bm{B}^F(k)$ as follows:

\begin{equation}\label{eq:but_factor}
\begin{aligned}
    \bm{B}^F(k)=\begin{bmatrix}
    \text{diag}(\bm{d}_1) & \text{diag}(\bm{d}_2) \\
    \text{diag}(\bm{d}_3) & \text{diag}(\bm{d}_4)
    \end{bmatrix}\in\mathbb{R}^{k\times k},
\end{aligned}
\end{equation}

where each $\bm{d}_i\in\mathbb{R}^{\frac{k}{2}}$ for $i = 1, 2, 3, 4$ are vectors. Next, assuming $d=2^N$, the $d$-dimensional \emph{butterfly matrix} $\bm{B}(d)\in\mathbb{R}^{d\times d}$ is recursively specified by

\begin{equation}
\begin{aligned}
    \!\!\bm{B}(d)=\tilde{\bm{B}}(d,d)\!\cdot\!\begin{bmatrix}
    \bm{B}_1(\frac{d}{2})\!\!\!\!\! & \bm{0} \\
    \bm{0}\!\!\!\!\! &\bm{B}_2(\frac{d}{2})
    \end{bmatrix}= \tilde{\bm{B}}(d,d)\tilde{\bm{B}}(d,\frac{d}{2})\cdots\tilde{\bm{B}}(d,2),
\end{aligned}
\end{equation}

where $\bm{B}_1(\frac{d}{2})$ and $\bm{B}_2(\frac{d}{2})$ denote two butterfly matrices of dimension $\frac{d}{2}$. We introduce the \emph{butterfly component} as $\thickmuskip=2mu \medmuskip=2mu \tilde{\bm{B}}(d,k)=\text{diag}(\bm{B}^F_1(k),\cdots,\bm{B}^F_{d/k}(k))$, representing a block-diagonal matrix with size $d\times d$ whose individual blocks are of dimension $k$. Here, each $\bm{B}^F_i(k)$ is a butterfly factor, as defined in Equation~\ref{eq:but_factor}. This setup allows us to construct an orthogonal matrix parameterization leveraging butterfly matrices. To ensure orthogonality, it suffices for all butterfly matrix factors $\tilde{\bm{B}}(d,k)$ composing $\bm{B}(d)$ to be orthogonal matrices. Focusing first on $\tilde{\bm{B}}(d,2)$—a block-diagonal matrix with blocks of size $2$—orthogonality is straightforward to enforce by parameterizing each $2\times 2$ block with the Cayley transform (equivalent to a $2$-dimensional rotation), consistent with \cite{liu2021orthogonal,qiu2023controlling}. Since the sparsity pattern of every butterfly component is a permutation of the pattern in $\tilde{\bm{B}}(d,2)$, all such components can be similarly parameterized to be orthogonal. This yields a compact and effective way to model orthogonal matrices using a collection of $2\times 2$ orthogonal matrices. Adopting the approach from \cite{chen2022pixelated}, we extend the definition to a block butterfly component $\tilde{\bm{B}}^b(d,k)$, where each element $\bm{d}_i$ is now replaced by a $b\times b$ matrix. To enforce orthogonality in the block butterfly component $\tilde{\bm{B}}^b(d,2)$, we ensure that every $2b\times 2b$ block is parameterized as an orthogonal matrix. For other butterfly components $\tilde{\bm{B}}^b(d,k)$ with $k>2$, the non-zero structure is given by a block permutation of the non-zero structure in $\tilde{\bm{B}}^b(d,2)$, thus permitting similar orthogonal parameterizations. Integrating these pieces, the forward computation in BOFT is given by

\begin{equation}
    \begin{aligned}\nonumber
    \bm{z}=\big(\bm{R}(m,b)\cdot\bm{W}^0\big)^\top\bm{x},~~~~\text{s.t.}~\left\{\bm{R}(m,b)=\prod_{i=1}^m \tilde{\bm{B}}^b_{(d,i)}~~\&~~\underbrace{\big(\tilde{\bm{B}}^b_{(d,j)}\big)^\top\tilde{\bm{B}}^b_{(d,j)}=\tilde{\bm{B}}^b_{(d,j)}\big(\tilde{\bm{B}}^b_{(d,j)}\big)^\top\!=\bm{I}_d}_{\forall j\in [1, m]}\right\},
    \end{aligned}
\end{equation}

Here, we use the notation $\tilde{\bm{B}}^b(d,2^{m - i+1})$ as a shorthand for $\tilde{\bm{B}}^b_{(d,i)}$, and $\bm{I}_d$ represents the $d$-dimensional identity matrix. The matrix $\bm{R}(m,b)\in\mathbb{R}^{d\times d}$ is an orthogonal matrix formed by multiplying several orthogonal butterfly components. For practical purposes, we refer to the BOFT parameterization employing $\bm{R}(m,\frac{b}{2})$ as BOFT($m$,$b$), where $b\geq 2$ must hold. In the special case where $\thickmuskip=2mu \medmuskip=2mu m=1$, BOFT($1$,$b$) becomes equivalent to the block-diagonal OFT~\cite{qiu2023controlling} with block size $b$. If $b=d$, i.e., BOFT($1$,$d$), the construction matches the original OFT~\cite{qiu2023controlling} with a completely unconstrained orthogonal matrix. Employing BOFT($\log \frac{2d}{b}$,$b$) selects the block butterfly matrix $\bm{B}^{\frac{b}{2}}(d)$ as $\bm{R}$, producing a dense orthogonal transformation. Broadly, BOFT($m$,$b$) uses a total of $\frac{1}{2}(b-1)dm$ learnable parameters for adaptation of a linear transformation with dimensions $d\times n$. When leveraging the standard butterfly matrix with $m=\log d$ and $b=2$, the parameter count scales as $\mathcal{O}(d\log d)$. On the other hand, the classical OFT based on a fully dense orthogonal matrix scales as $\mathcal{O}(d^2)$, while the block-diagonal OFT with $r$ blocks requires $\mathcal{O}(bd)$. As a consequence, producing a dense orthogonal matrix in the original OFT necessitates a block size of $b=d$, but BOFT is flexible and can construct such matrices with any $b$.

\textbf{Identity Initialization in BOFT}. When finetuning, it is standard to initialize with weights identical to those in the pretrained model, thereby ensuring minimal drift between the pretrained and adapted models. For instance, LoRA adopts an initialization of zeros for the additional low-rank weights. In BOFT, we similarly initialize all butterfly components as identity matrices (i.e., the associated skew-symmetric matrix for the Cayley transform is set to zeros).

\textbf{Multiplicative Dropout Mechanism in BOFT}. LoRA~\cite{hulora2022} incorporates a Dropout operation within its low-rank weight adaptation in order to combat overfitting. While conventional Dropout~\cite{srivastava2014dropout} is suitable for LoRA, it is incompatible with BOFT due to its multiplicative nature of weight updates. To resolve this, we introduce a new multiplicative Dropout specifically tailored for BOFT. Noting that $\bm{R}(m,b)$ is composed as a sequence of $m$ orthogonal butterfly components that can be decomposed into $2b\times 2b$ block-diagonal orthogonal submatrices, this Dropout scheme involves randomly selecting a proportion $p_1$ of butterfly components and a fraction $p_2$ of the diagonal blocks within each, and subsequently replacing those chosen blocks by identity matrices.

\section{Intriguing Insights and Discussions}

\textbf{Expressivity of BOFT}. While the butterfly structure, in combination with permutation operations, has been shown to exactly recover various canonical fast linear transforms~\cite{dao2019learning,dao2019kaleidoscope} (\emph{e.g.}, the fast Fourier transform, Hadamard transform), the extent to which our orthogonal butterfly matrix is capable of approximating arbitrary orthogonal matrices has yet to be thoroughly explored. To investigate this, we perform simulations wherein we approximate a randomly sampled dense orthogonal matrix~\cite{anderson1987generation} of dimensions $512\times 512$. The outcomes, presented in Figure~\ref{fig:para_eff}, are averaged over 10 distinct random seeds. The vertical axis indicates the approximation error, while the horizontal axis reflects the number of effective trainable parameters. Each curve of the same color corresponds to BOFT using a particular block size, with the leftmost marker representing the approximation error of BOFT($1$,$b$) (\emph{i.e.}, a traditional block-diagonal OFT with block size $b$). Overall, BOFT consistently provides improved parameter efficiency relative to OFT. For instance, BOFT($9$,$2$) demonstrates higher expressivity than BOFT($1$,$16$), yet involves significantly fewer parameters. Likewise, configurations with smaller $b$ and larger $m$ for BOFT typically achieve greater parameter efficiency; BOFT($6$,$4$) employs far fewer parameters while attaining comparable approximation error to BOFT($2$,$16$). More generally, the butterfly matrix encompasses a highly structured subspace within the full orthogonal group—distinct from the subspace induced by block-diagonal matrices—rendering BOFT strictly more expressive than OFT at a fixed block size.

\begin{theorem}[Expressivity of BOFT]\label{thm:exp_boft}
    BOFT surpasses OFT in expressivity given an identical block size. It is possible to approximate every orthogonal matrix of dimension $d$ by composing butterfly matrices of the form $\bm{B}_{d-1,1}(d)\bm{B}^\top_{d-1,2}(d)\cdots\bm{B}_{1,1}(d)\bm{B}^\top_{1,2}(d)$, where $\bm{B}_{i,j}(d)$, for all $i, j$, are butterfly matrices.
\end{theorem}

Building upon Theorem~\ref{thm:exp_boft}, we can further generalize BOFT so that the target orthogonal matrix is expressed as $\thickmuskip=2mu \medmuskip=2mu \bm{R}^G(m_1,b_1,m_2,b_2,l)=\bm{R}_{l,1}(m_1,b_1)\bm{R}_{l,2}^T(m_2,b_2)\cdots\bm{R}_{1,1}(m_1,b_1)\bm{R}_{1,2}^T(m_2,b_2)$, where $\bm{R}_{i,j}^T(m,b)$ refers to the orthogonal matrix employed within BOFT. Notably, when parameters are set to $m_1=m_2=\log d$, $b_1=b_2=2$, and $l=d-1$, the generalized matrix $\bm{R}^G(m_1,b_1,m_2,b_2,l)$ is able to span the entire orthogonal group. This recursive matrix construction is also known as the kaleidoscope hierarchy~\cite{dao2019kaleidoscope}. Nevertheless, it is important to recognize that increased expressiveness does not guarantee enhanced finetuning performance—full finetuning, which is maximally expressive, may in practice deliver suboptimal results. The essential challenge lies in balancing expressivity with regularization, as this determines the model’s generalization capacity during finetuning. BOFT expands the finetuning search space in a principled, structured manner, thereby allowing us to better negotiate the trade-off between expressiveness and inductive regularity.

\textbf{Spectral properties}. Compared to LoRA, orthogonal finetuning consistently achieves superior spectral characteristics, owing to its exact retention of the spectral norm of the pretrained weight matrix $\bm{W}^0$. This preservation is evident when considering the singular value decomposition: $\bm{W}^0=\bm{U}\bm{\Sigma}\bm{V}^\top$, with $\bm{U}$ and $\bm{V}$ as orthogonal matrices and $\bm{\Sigma}$ containing the singular values along its diagonal. Both OFT and BOFT apply a left multiplication by an orthogonal matrix $\bm{R}$, resulting in updated weights $\bm{R}\bm{U}\bm{\Sigma}\bm{V}^\top$. This operation maintains the largest singular value—hence, the spectral norm—of $\bm{W}^0$ intact. Such invariance has been demonstrated to markedly enhance both the stability of training processes and the generalization capacity of models~\cite{miyato2018spectral,yoshida2017spectral}. Additional mathematical discussions related to these properties are provided in Appendix~\ref{app:sn_property}.

\textbf{Orthogonal finetuning as learning bilinear similarity}. BOFT may also be interpreted as a mechanism for learning bilinear similarities, as in $\bm{w}^0_i\bm{R}\bm{x}$, where $\bm{w}^0_i$ denotes the $i$-th column (or neuron) of the original weight matrix $\bm{W}^0$. In this context, BOFT frames the learning process as identifying a bilinear similarity transformation matrix $\bm{R}$, subject to a stringent constraint—orthogonality—which relates closely to classic ideas from distance metric learning~\cite{xing2002distance} and the study of bilinear forms~\cite{roman2005advanced}.

\textbf{Inductive bias and generalization in BOFT}. In the BOFT framework, the matrix $\bm{R}(m,b)$ typically selects from a well-structured subset within the orthogonal group, thereby introducing an explicit constraint on the space of solutions. This leads to an inherent inductive bias. We posit that the inductive bias imparted through butterfly factorization in BOFT fosters better generalization, since it captures the underlying structure shared by numerous foundational linear transformations, such as the discrete Fourier transform, discrete sine/cosine transform, and Hadamard transform~\cite{dao2019kaleidoscope}. Furthermore, the use of sparse factorization within BOFT is also expected to confer additional, potentially beneficial, implicit biases~\cite{gunasekar2017implicit,li2018algorithmic,liu2019neural}.

\textbf{Relation to butterfly-based sparse training}. Numerous prior studies~\cite{dao2019learning,dao2019kaleidoscope,chen2022pixelated,dao2022monarch} have investigated sparse training utilizing the butterfly parameterization. These works generally employ direct reparameterization of weight matrices via butterfly structures and train the corresponding neural networks from the ground up. The approach in \cite{dao2022monarch} deviates slightly by applying finetuning to pretrained weights; this is done by first projecting weights onto a form of butterfly matrices and then optimizing the resultant components to adapt to specific tasks. In contrast, BOFT introduces a novel finetuning procedure wherein weights are modified using weight matrices that are shared across layers.

\input{rebuttal-results}

\section{Concluding Remarks and Limitations}

In this work, we present Orthogonal Butterfly, a broadly applicable and parameter-efficient finetuning approach for foundation models leveraging the butterfly matrix structure. The principal contribution to enhancing parameter efficiency is the parameterization of a dense orthogonal matrix as the product of several sparse orthogonal matrices. To systematically identify viable matrix decompositions, we introduce a graphical information transmission perspective. This framework reveals that butterfly matrices are particularly well suited for achieving sparse orthogonal factorizations aligned with our objectives. Through empirical study, we establish that BOFT achieves strong finetuning performance on a range of tasks, including large language models, extensive vision systems, and text-to-image generation. Our experimental results further highlight BOFT's capability as a versatile approach to finetuning.

Nevertheless, BOFT is not without shortcomings. The construction of the final orthogonal matrix in BOFT as a product of multiple orthogonal matrices leads to a somewhat increased training runtime relative to OFT, posing an open challenge for future acceleration. It is worth noting, however, that post-finetuning, the orthogonal matrices obtained by BOFT can be merged directly into the pretrained parameters, incurring no extra inference cost. Additionally, it remains unclear whether the butterfly network represents the most effective mechanism for information transmission. Our proposed information transmission view facilitates exploration beyond traditional machine learning by drawing on insights from computer networking—a field that rigorously examines the effectiveness of network topologies in communication. We anticipate that exploiting alternative network architectures may offer further improvements for dense orthogonal matrix composition.

\end{document}